\section{Общая схема метода и программного конвейера}

Методология работы состоит из двух связанных частей. Первая часть --- программный конвейер построения данных, генерации кандидатов и автоматической оценки. Вторая часть --- обучающая цель MRVF, то есть многореференсное обучение без внешнего верификатора, использующее найденные эталонные ответы как слабый сигнал для обучения траекторий рассуждения.

На высоком уровне система строит слабо размеченный набор данных с несколькими эталонными ответами из сырых корпусов шуток. Сначала шутки объединяются в единый корпус и очищаются от точных дублей и почти совпадающих текстов. Затем для каждой шутки строится плотное векторное представление, из текста извлекаются ключевые слова, а группы ключевых слов используются как слабые запросы для поиска семантически близких шуток. Полученные множества близких шуток образуют наблюдаемые наборы эталонных ответов \(Y^{*}(x)\).

Эти эталонные множества используются в двух ролях. Во-первых, они задают контрольные запросы для генерации кандидатов и последующей автоматической оценки. Во-вторых, они входят непосредственно в обучающую функцию MRVF: модель получает запрос \(x\), порождает траекторию рассуждения \(z\), после чего эталонные ответы из \(Y^{*}(x)\) оцениваются с помощью принудительного учительского декодирования. Таким образом, подготовка данных не является вспомогательным этапом: она определяет наблюдаемую структуру, относительно которой формулируется обучающая цель.

\section{Построение набора данных}

Набор данных объединяется из двух публичных коллекций шуток: Short Jokes и rJokes \citep{weller2020rjokes}. В репозитории этот этап реализован в \texttt{src/pipelines/jokes.py}. Конвейер скачивает сырые файлы из зафиксированных версий источников, преобразует их в формат Parquet и сохраняет метаданные источника для каждой шутки. Для каждой строки сохраняются поля \texttt{id}, \texttt{text}, \texttt{source\_name}, \texttt{source\_filename} и \texttt{source\_id}. После отдельной предобработки двух корпусов они объединяются и получают новый глобальный целочисленный идентификатор.

Исходная коллекция содержит приблизительно \(664\,000\) текстов шуток. После удаления точных дублей и почти совпадающих текстов размер текущего корпуса составляет приблизительно \(570\,000\) шуток. Дедупликация важна не только как техническая очистка данных. В многореференсной постановке повторяющаяся шутка фактически получает больший вес в наблюдаемом множестве эталонных ответов, что может искусственно усиливать её вклад в обучающую цель. Поэтому удаление дублей является частью определения эмпирической меры над корпусом шуток.

\subsection{Удаление дублей и почти совпадающих шуток}

Для корпуса юмора полная агрессивная нормализация нежелательна: пунктуация, порядок слов и точная формулировка могут влиять на комический эффект. Поэтому очистка выполняется в несколько этапов. Сначала удаляются пустые записи и точные дубли после нормализации строки. Пусть \(n(y)\) обозначает нормализованный текст шутки \(y\), полученный после приведения Unicode к канонической форме, приведения к нижнему регистру и удаления неинформативных символов. Две шутки считаются точными дублями, если
\[
    n(y_i)=n(y_j).
\]

Следующий уровень удаляет тексты, совпадающие по множеству значимых токенов. Пусть \(T(y)\) --- множество токенов шутки после простой токенизации и удаления служебных символов. Тогда токенный дубль определяется условием
\[
    T(y_i)=T(y_j),
\]
при условии, что число уникальных токенов достаточно велико. Этот дополнительный порог нужен, чтобы не удалять короткие тексты, которые случайно совпадают по одному или двум словам.

Наконец, для почти совпадающих шуток используются меры сходства по токенам и символьным фрагментам. Токенное сходство Жаккара определяется как
\[
    J_{\mathrm{tok}}(y_i,y_j)
    =
    \frac{|T(y_i)\cap T(y_j)|}{|T(y_i)\cup T(y_j)|}.
\]
Аналогично вычисляется сходство Жаккара по символьным шинглам \(J_{\mathrm{char}}\). Дополнительно используется нормированное сходство редактирования \(R_{\mathrm{edit}}\), оценивающее близость строк с учётом порядка символов. Пара шуток считается почти дублем, если соответствующие значения превышают заданные пороги:
\[
    J_{\mathrm{tok}}(y_i,y_j)\ge \tau_{\mathrm{tok}},
    \qquad
    J_{\mathrm{char}}(y_i,y_j)\ge \tau_{\mathrm{char}},
    \qquad
    R_{\mathrm{edit}}(y_i,y_j)\ge \tau_{\mathrm{edit}}.
\]

Такой подход сохраняет консервативность очистки: корпус не сводится к абстрактным семантическим кластерам, но очевидные повторения и переформулировки одной и той же шутки удаляются. Это особенно важно для дальнейшего поиска эталонных ответов: если дубликаты оставить, то один и тот же юмористический ход будет многократно попадать в \(Y^{*}(x)\) и завышать свою роль при обучении.

\section{Плотный индекс векторных представлений}

Следующий этап кодирует каждую шутку в плотном семантическом векторном пространстве. Он реализован в \texttt{src/pipelines/embeddings.py}. Текущая конфигурация использует модель семейства Qwen Embedding и хранит для каждой шутки вектор фиксированной размерности. Векторные представления записываются во фрагменты Parquet и хранятся как пары \((i,e_i)\), где \(i\) --- глобальный идентификатор шутки, а \(e_i\in\mathbb{R}^{d}\) --- соответствующее представление.

Роль этого этапа двоякая. Во-первых, векторные представления используются для оценки релевантности ключевых слов-кандидатов. Во-вторых, они задают индекс поиска, с помощью которого для слабого запроса \(x\) извлекаются семантически близкие шутки, образующие наблюдаемое множество эталонных ответов \(Y^{*}(x)\).

Иными словами, пространство векторных представлений является общим представлением, от которого зависят и извлечение ключевых слов, и поиск эталонных ответов. Это важная инженерная особенность конвейера: качество многореференсной разметки зависит не от ручной разметки запросов, а от того, насколько хорошо векторное пространство группирует тематически и семантически близкие шутки.

\section{Извлечение ключевых слов как поиск слабых запросов}

Предполагается, что каждая шутка имеет набор слабых запросов, на которые она правдоподобно отвечает. Например, запрос ``write a joke about a chicken'' можно рассматривать как допустимый запрос для классической шутки ``why did the chicken cross the road''. Этап извлечения ключевых слов направлен на построение таких слабых запросов на основе текста самой шутки. Он реализован в \texttt{src/pipelines/keywords.py}.

Для каждой шутки \(y_i\) процесс сначала генерирует пул n-грамм-кандидатов с помощью анализатора \texttt{CountVectorizer}. В конфигурации по умолчанию этот анализатор использует униграммы, удаляет английские стоп-слова и оставляет ограниченное число терминов-кандидатов на шутку. Пусть полученный набор кандидатов равен
\[
    C(y_i)=\{c_{i1},\dots,c_{im}\}.
\]
Затем каждый термин-кандидат кодируется той же моделью векторных представлений, что и шутки.

Если \(e(y_i)\) обозначает векторное представление шутки, а \(e(c)\) --- векторное представление ключевого слова-кандидата, то релевантность кандидата оценивается через косинусное сходство:
\[
    s(c,y_i)
    =
    \frac{e(c)^\top e(y_i)}
    {\lVert e(c)\rVert\,\lVert e(y_i)\rVert}.
\]
Выбор только наиболее похожих кандидатов часто давал бы избыточные ключевые слова. Например, несколько близких n-грамм могут описывать один и тот же объект или одну и ту же тему.

Для уменьшения избыточности применяется maximal marginal relevance (MMR), то есть процедура выбора кандидатов с балансом между релевантностью и разнообразием. На каждом шаге выбирается кандидат, который одновременно хорошо соответствует исходной шутке и не слишком похож на уже выбранные ключевые слова. В результате для шутки получается компактное представление
\[
    K(y_i)=\{k_{i1},k_{i2},k_{i3}\},
\]
которое может содержать меньше трёх элементов, если пул кандидатов мал.

Методологически эти ключевые слова не следует интерпретировать как полностью восстановленные исходные запросы. Это более слабые объекты: короткие n-граммы, которые фиксируют тематические или лексические якоря шутки. Такое приближение не решает задачу восстановления запроса по шутке полностью, но даёт удобное представление, из которого можно строить контролируемые запросы для генерации.

\section{Построение эталонных ответов через группировку ключевых слов и поиск}

\subsection{Группировка ключевых слов и формирование запросов}

Этап построения эталонных ответов реализован в \texttt{src/pipelines/references.py}. Для каждого набора ключевых слов \(K(y_i)\) процесс перечисляет все уникальные комбинации ключевых слов, размер которых лежит между \texttt{min\_keywords} и \texttt{max\_keywords}. В текущей постановке используются однословные и двухсловные группы.

Если
\[
    K(y_i)=\{k_1,k_2,k_3\},
\]
то сгенерированные группы имеют вид
\[
    \{k_1\},\ \{k_2\},\ \{k_3\},\ \{k_1,k_2\},\ \{k_1,k_3\},\ \{k_2,k_3\}.
\]
Каждая группа преобразуется в запрос на естественном языке с помощью шаблона
\begin{quote}
    \ttfamily Write a joke using the following keyword(s): \{keywords\}
\end{quote}
Этот запрос служит двум целям. Во-первых, это текстовая форма, которая затем подаётся генеративным моделям. Во-вторых, этот же объект кодируется как поисковый запрос для извлечения эталонных ответов из корпуса.

Запросы с одним ключевым словом являются широкими, тематическими и часто допускают много возможных шуток. Запросы с двумя ключевыми словами более ограничены и обычно извлекают более точные кандидаты. Текущая реализация останавливается на двух ключевых словах, потому что более крупные группы быстро становятся слишком специфичными и уменьшают число правдоподобных различных эталонных ответов.

\subsection{Поиск эталонных ответов с помощью индекса векторных представлений}

После формирования строк запросов процесс кодирует их как поисковые запросы и извлекает близкие шутки из индекса FAISS, построенного по полной коллекции векторных представлений шуток. Перед индексацией и поиском векторы нормализуются, так что скалярное произведение соответствует косинусному сходству.

Для сформированного запроса \(x\) пусть \(q(x)\) --- его нормализованное векторное представление, а \(e_j\) --- нормализованное векторное представление шутки \(y_j\). Тогда оценка поиска равна
\[
    s(x,y_j)=q(x)^\top e_j.
\]
Процедура поиска извлекает больше кандидатов, чем требуется в финальном наборе, фильтрует те, чьё сходство ниже заданного порога, и оставляет лучшие оставшиеся объекты. Соответствующие тексты используются как наблюдаемый набор эталонных ответов:
\[
    Y^{*}(x)=\{y_1,\dots,y_n\},
    \qquad
    n\le 5.
\]

Важно, что \(Y^{*}(x)\) не является полным множеством всех хороших шуток для запроса \(x\). Это конечное множество найденных корпусных примеров, близких к запросу в выбранном векторном пространстве. Поэтому дальнейшая обучающая цель должна интерпретироваться как оптимизация по неполному наблюдаемому множеству эталонов, а не как оптимизация истинной человеческой смешности.

Этот этап реализует центральную идею работы: запрос связывается не с одним истинным ответом, а с несколькими шутками-кандидатами, которые могут служить эталонными ответами. Такая постановка лучше соответствует природе юмора, где один и тот же набор ключевых слов допускает множество разных удачных панчлайнов.

\subsection{Дедупликация эталонных запросов и разбиение выборки}

После поиска итоговый набор данных содержит строки вида
\[
    (x_i, Y^{*}(x_i), S(x_i)),
\]
где \(x_i\) --- запрос на основе ключевых слов, \(Y^{*}(x_i)\) --- найденный набор эталонных ответов, а \(S(x_i)\) --- вектор оценок поиска. Поскольку разные исходные шутки могут порождать одинаковые группы ключевых слов, набор эталонных запросов дополнительно дедуплицируется по полю \texttt{keywords}. При этом сохраняется первое вхождение, а идентификаторы строк назначаются заново.

Дедуплицированный набор затем разбивается на обучающую, валидационную и тестовую части. Разбиение выполняется случайно с фиксированным зерном, что обеспечивает воспроизводимость экспериментов. Отдельно сохраняются полный набор данных и каждая часть выборки. Такая организация позволяет использовать один и тот же набор эталонных запросов для обучения, генерации кандидатов и последующей автоматической оценки.

\section{Генерация кандидатов}

Этап генерации кандидатов поддерживает сравнение нескольких вариантов модели на одном и том же наборе запросов. Он реализован в \texttt{src/pipelines/candidates.py}. Для выбранной модели и части набора данных процесс проходит по запросам и вызывает интерфейс генерации ответов, совместимый с OpenAI API.

Используется тот же общий шаблон запроса, что и при построении эталонных ответов:
\begin{quote}
    \ttfamily Write a joke using the following keyword(s): \{keywords\}
\end{quote}
Конкретные экспериментальные запуски могут уточнять этот шаблон, например явно требовать вернуть только финальную шутку без объяснений. Такое уточнение важно, поскольку в задаче юмора формат ответа является частью качества: длинное объяснение или повторение инструкции может ухудшать оценку даже при наличии удачной идеи.

Схема выхода содержит четыре основных поля: идентификатор строки, группу ключевых слов, название модели и сгенерированный текст. Файлы с кандидатами сохраняются в структуре каталогов по модели и части выборки. Такая организация важна для последующей оценки, поскольку позволяет выравнивать ответы нескольких моделей по одному базовому идентификатору и запросу.

Методологически этот этап следует понимать как слой построения контрольных артефактов. Он создаёт ответы нескольких вариантов модели на одни и те же запросы: базовых моделей, моделей в режиме явного рассуждения и моделей после MRVF-обучения. Благодаря общему идентификатору строки ответы разных моделей можно выравнивать по одному запросу и сравнивать попарно.

\section{Оценка кандидатов}

\subsection{Попарное сравнение}

Автоматическая оценка реализована в \texttt{src/pipelines/evaluation.py}. Процесс сначала объединяет наборы кандидатов от нескольких моделей. Затем ответы группируются по идентификатору эталонного запроса и названию модели. Для каждого запроса строятся все попарные сравнения между моделями, которые сгенерировали непустые ответы.

Порядок левого и правого ответа рандомизируется с фиксированным зерном случайности, чтобы уменьшить позиционное смещение модели-оценщика. Оценочный запрос содержит исходный запрос на шутку, левый и правый кандидатные ответы. Системная инструкция просит модель-оценщик выбрать, какая шутка смешнее для широкой аудитории, избегать ничьих и возвращать строгий JSON-ответ. В текущем протоколе модель-оценщик используется при температуре \(0\).

Такой слой оценки измеряет только сравнительное качество ответов в рамках выбранной модели-оценщика. Он не заменяет человеческую оценку и не реализует многомерную рубрику, включающую новизну, релевантность, краткость или социальную уместность. Тем не менее попарное сравнение является удобным промежуточным инструментом, поскольку относительный выбор между двумя шутками обычно проще и устойчивее абсолютной числовой оценки.

\subsection{Агрегация таблицы результатов}

Попарные суждения агрегируются в рейтинг с числом побед, числом поражений, числом сравнений, долей побед и оценками Брэдли--Терри. Пусть \(\mathcal{M}\) обозначает множество оцениваемых моделей. Для каждой пары моделей \((m_a,m_b)\) и каждого запроса модель-оценщик выбирает победителя.

Модель Брэдли--Терри сопоставляет каждой модели \(m\) скрытый параметр качества \(q_m\). Вероятность победы модели \(m_a\) над моделью \(m_b\) задаётся как
\[
    P(m_a \succ m_b)
    =
    \frac{\exp(q_{m_a})}
    {\exp(q_{m_a})+\exp(q_{m_b})}.
\]
В таблице результатов используется оценённое значение \(q_m\), обозначаемое как \texttt{bt\_score}. Такая агрегация позволяет получить единую относительную шкалу качества из набора попарных сравнений, но её следует интерпретировать только относительно выбранного набора моделей, запросов и модели-оценщика.

Таким образом, программный конвейер решает две задачи. С инженерной стороны он создаёт воспроизводимые артефакты: очищенный корпус, векторные представления, ключевые слова, эталонные множества, кандидатные ответы и таблицы автоматической оценки. С методологической стороны он задаёт наблюдаемую структуру \(Y^{*}(x)\), без которой невозможно сформулировать MRVF-обучение. Дальнейшие разделы описывают, как эти множества эталонных ответов входят в обучающую цель.

\section{Теоретическая рамка обучения MRVF}

Построенный выше конвейер связывает каждый запрос \(x\) с наблюдаемым множеством эталонных ответов \(Y^{*}(x)\). Теперь опишем, как такая многореференсная разметка входит в обучающую цель. Центральная идея состоит в том, что генерация юмора плохо согласуется с предположением об одном правильном ответе. Слабый запрос может допускать много приемлемых шуток, различающихся завязкой, механизмом панчлайна, стилем и уровнем абсурдности. Поэтому обучающая цель должна учитывать не один эталонный ответ, а множество наблюдаемых эталонов.

\subsection{Обозначения}

Пусть \(\mathcal{X}\), \(\mathcal{Z}\) и \(\mathcal{Y}\) обозначают пространства запросов, траекторий рассуждения и итоговых ответов соответственно. Для запроса \(x\in\mathcal{X}\):
\begin{itemize}
    \item \(z\in\mathcal{Z}\) обозначает траекторию рассуждения;
    \item \(y\in\mathcal{Y}\) обозначает итоговую сгенерированную шутку;
    \item \(y^{*}\in\mathcal{Y}\) обозначает один эталонный ответ;
    \item \(Y^{*}(x)\subset\mathcal{Y}\) обозначает наблюдаемое множество эталонных ответов, связанное с запросом \(x\);
    \item \(\pi_{\theta}\) обозначает стратегию, то есть условную вероятностную модель над текстом, параметризованную \(\theta\).
\end{itemize}

Совместное распределение по траектории рассуждения и итоговому ответу записывается как
\[
    \pi_{\theta}(z,y\mid x)
    =
    \pi_{\theta}(z\mid x)\,
    \pi_{\theta}(y\mid x,z).
\]
Когда модель явно выводит токены рассуждения, \(z\) можно отождествить с текстом внутри специального сегмента рассуждения, а \(y\) --- с финальной шуткой. Если явное рассуждение отключено, ту же факторизацию можно понимать как концептуальное разложение поведения модели на выбор промежуточного состояния и последующую генерацию ответа.

При обсуждении неполноты будет полезно представлять более крупное скрытое множество
\[
    Y_{\mathrm{true}}(x)\supseteq Y^{*}(x),
\]
содержащее все действительно приемлемые шутки для запроса \(x\). Это полное множество никогда не наблюдается на практике, но оно помогает формально описать ограничение многореференсной разметки.

\subsection{Обучение с верифицируемыми наградами}

Обучение с подкреплением оптимизирует стратегию, максимизируя ожидаемую награду. В задачах с верифицируемым ответом награда может быть назначена внешней процедурой, не зависящей от субъективного человеческого суждения. Например, программу можно выполнить на наборе тестов, а итоговый ответ в математической задаче можно сравнить с каноническим решением. Именно такая структура лежит в основе ряда современных работ по обучению рассуждающих моделей \citep{DeepSeek2025,Lightman2023}.

Сначала рассмотрим стандартную постановку с одним эталонным ответом. Для каждого запроса \(x\) существует проверяемая цель \(y^{*}\), а награда задаётся как
\[
    R(y,y^{*})=\mathbb{I}\{y=y^{*}\},
\]
где \(\mathbb{I}\{\cdot\}\) --- индикаторная функция. Тогда целевая функция имеет вид
\[
    J_{\mathrm{RLVR}}(\theta\mid x,y^{*})
    =
    \mathbb{E}_{z\sim\pi_{\theta}(z\mid x)}
    \mathbb{E}_{y\sim\pi_{\theta}(\cdot\mid x,z)}
    \big[
        R(y,y^{*})
        \big].
\]
Эквивалентно, в развернутой форме суммы:
\[
    J_{\mathrm{RLVR}}(\theta\mid x,y^{*})
    =
    \sum_{z\in\mathcal{Z}}
    \sum_{y\in\mathcal{Y}}
    \pi_{\theta}(z\mid x)
    \pi_{\theta}(y\mid x,z)
    R(y,y^{*}).
\]

Для оптимизации этого математического ожидания используется тождество функции оценки, также известное как тождество REINFORCE:
\[
    \nabla_{\theta}
    \mathbb{E}_{u\sim p_{\theta}(u)}
        [f(u)]
    =
    \mathbb{E}_{u\sim p_{\theta}(u)}
    \big[
        f(u)\nabla_{\theta}\log p_{\theta}(u)
        \big].
\]
Это тождество следует напрямую из соотношения
\[
    \nabla_{\theta}p_{\theta}(u)
    =
    p_{\theta}(u)
    \nabla_{\theta}\log p_{\theta}(u).
\]

Применяя его к совместному распределению \((z,y)\), получаем
\begin{align*}
    \nabla_{\theta}J_{\mathrm{RLVR}}
     & =
    \nabla_{\theta}
    \sum_{z,y}
    \pi_{\theta}(z,y\mid x)
    R(y,y^{*})
    \\
     & =
    \sum_{z,y}
    R(y,y^{*})
    \nabla_{\theta}
    \pi_{\theta}(z,y\mid x)
    \\
     & =
    \sum_{z,y}
    \pi_{\theta}(z,y\mid x)
    R(y,y^{*})
    \nabla_{\theta}
    \log\pi_{\theta}(z,y\mid x)
    \\
     & =
    \mathbb{E}_{z,y}
    \Big[
        R(y,y^{*})
        \nabla_{\theta}
        \log\pi_{\theta}(z,y\mid x)
        \Big].
\end{align*}
Используя факторизацию совместной стратегии,
\[
    \log\pi_{\theta}(z,y\mid x)
    =
    \log\pi_{\theta}(z\mid x)
    +
    \log\pi_{\theta}(y\mid x,z),
\]
получаем
\[
    \nabla_{\theta}J_{\mathrm{RLVR}}
    =
    \mathbb{E}_{z,y}
    \Big[
        R(y,y^{*})
        \big(
        \nabla_{\theta}\log\pi_{\theta}(z\mid x)
        +
        \nabla_{\theta}\log\pi_{\theta}(y\mid x,z)
        \big)
        \Big].
\]

Это выражение имеет простую интерпретацию. Выбранная пара \((z,y)\) получает награду только тогда, когда верификатор объявляет \(y\) правильным. Когда это происходит, обновление увеличивает и вероятность траектории рассуждения, приведшей к ответу, и вероятность самого ответа при этой траектории.

\subsection{VeriFree}

Юмор не предоставляет такого внешнего верификатора, и в более общем виде многие задачи открытой генерации не имеют детерминированного теста правильности. VeriFree закрывает этот разрыв, замечая, что в случае одного эталонного ответа ожидаемую награду точного совпадения можно переписать как условную вероятность \citep{Zhou2025}. Зафиксируем \(x\) и \(z\). Тогда
\begin{align*}
    \mathbb{E}_{y\sim\pi_{\theta}(\cdot\mid x,z)}
    \big[
        \mathbb{I}\{y=y^{*}\}
        \big]
     & =
    \sum_{y\in\mathcal{Y}}
    \pi_{\theta}(y\mid x,z)
    \mathbb{I}\{y=y^{*}\}
    \\
     & =
    \pi_{\theta}(y^{*}\mid x,z).
\end{align*}
Подставляя это тождество обратно в целевую функцию RLVR, получаем
\[
    J_{\mathrm{RLVR}}(\theta\mid x,y^{*})
    =
    \mathbb{E}_{z\sim\pi_{\theta}(z\mid x)}
    \big[
        \pi_{\theta}(y^{*}\mid x,z)
        \big].
\]
Это мотивирует награду без внешнего верификатора
\[
    r_{\theta}(x,z,y^{*})
    :=
    \pi_{\theta}(y^{*}\mid x,z)
\]
и соответствующую цель
\[
    J_{\mathrm{VF}}(\theta\mid x,y^{*})
    :=
    \mathbb{E}_{z\sim\pi_{\theta}(z\mid x)}
    \big[
        r_{\theta}(x,z,y^{*})
        \big].
\]

Значимость такой переформулировки одновременно практическая и концептуальная. Для вычисления награды не требуется явное декодирование итогового ответа \(y\). Вероятность \(\pi_{\theta}(y^{*}\mid x,z)\) можно вычислить с помощью принудительного учительского декодирования на последовательность эталонного ответа, что удобно параллелизуется и убирает шум, связанный с выборкой финального ответа.

Чтобы получить градиент, запишем
\[
    J_{\mathrm{VF}}(\theta\mid x,y^{*})
    =
    \sum_{z\in\mathcal{Z}}
    \pi_{\theta}(z\mid x)
    r_{\theta}(x,z,y^{*}).
\]
Дифференцируя по правилу произведения, получаем
\begin{align*}
    \nabla_{\theta}J_{\mathrm{VF}}
     & =
    \sum_z
    \Big[
        \nabla_{\theta}\pi_{\theta}(z\mid x)
        r_{\theta}(x,z,y^{*})
        +
        \pi_{\theta}(z\mid x)
        \nabla_{\theta}r_{\theta}(x,z,y^{*})
        \Big]
    \\
     & =
    \sum_z
    \pi_{\theta}(z\mid x)
    \Big[
        r_{\theta}(x,z,y^{*})
        \nabla_{\theta}\log\pi_{\theta}(z\mid x)
        +
        \nabla_{\theta}r_{\theta}(x,z,y^{*})
        \Big].
\end{align*}
Так как \(r_{\theta}(x,z,y^{*})=\pi_{\theta}(y^{*}\mid x,z)\), имеем
\[
    \nabla_{\theta}r_{\theta}(x,z,y^{*})
    =
    r_{\theta}(x,z,y^{*})
    \nabla_{\theta}\log\pi_{\theta}(y^{*}\mid x,z).
\]
Следовательно,
\[
    \nabla_{\theta}J_{\mathrm{VF}}
    =
    \mathbb{E}_{z\sim\pi_{\theta}(z\mid x)}
    \Big[
        r_{\theta}(x,z,y^{*})
        \big(
        \nabla_{\theta}\log\pi_{\theta}(z\mid x)
        +
        \nabla_{\theta}\log\pi_{\theta}(y^{*}\mid x,z)
        \big)
        \Big].
\]

Эта декомпозиция важна для дальнейшей работы. Первое слагаемое обновляет распределение траекторий рассуждения: оно повышает вероятность тех \(z\), при которых сама модель назначает большую условную вероятность эталонному ответу. Второе слагаемое является прямым градиентом по эталонному ответу, вычисляемым с помощью принудительного учительского декодирования.

\subsection{Многореференсное обобщение VeriFree}

Предположение об одном эталонном ответе не подходит для генерации юмора. Запрос вида ``напиши шутку с ключевыми словами cat и dog'' не имеет единственного правильного продолжения: возможны разные завязки, разные панчлайны и разные стилистические решения. Поэтому награда должна задаваться не на одном эталонном ответе, а на наблюдаемом множестве допустимых ответов.

Пусть \(Y^{*}(x)\subset\mathcal{Y}\) обозначает наблюдаемое множество эталонных ответов для запроса \(x\). Определим многореференсную награду точного совпадения:
\[
    R(y,Y^{*})
    =
    \mathbb{I}\{y\in Y^{*}\}.
\]
Соответствующая целевая функция имеет вид
\[
    J_{\mathrm{MR}}(\theta\mid x,Y^{*})
    =
    \mathbb{E}_{z\sim\pi_{\theta}(z\mid x)}
    \mathbb{E}_{y\sim\pi_{\theta}(\cdot\mid x,z)}
    \big[
        \mathbb{I}\{y\in Y^{*}\}
        \big].
\]
Для фиксированных \(x\) и \(z\)
\begin{align*}
    \mathbb{E}_{y\sim\pi_{\theta}(\cdot\mid x,z)}
    \big[
        \mathbb{I}\{y\in Y^{*}\}
        \big]
     & =
    \sum_{y\in\mathcal{Y}}
    \pi_{\theta}(y\mid x,z)
    \mathbb{I}\{y\in Y^{*}\}
    \\
     & =
    \sum_{y\in Y^{*}}
    \pi_{\theta}(y\mid x,z).
\end{align*}

Это мотивирует многореференсную награду без внешнего верификатора:
\[
    r_{\theta}(x,z,Y^{*})
    :=
    \sum_{y\in Y^{*}}
    \pi_{\theta}(y\mid x,z).
\]
Величина \(r_{\theta}(x,z,Y^{*})\) является вероятностной массой, которую модель при траектории \(z\) назначает наблюдаемому множеству эталонных ответов. Итоговая идеальная цель:
\[
    J_{\mathrm{MRVF}}(\theta\mid x,Y^{*})
    =
    \mathbb{E}_{z\sim\pi_{\theta}(z\mid x)}
    \big[
        r_{\theta}(x,z,Y^{*})
        \big].
\]

На уровне скалярной целевой функции MRVF играет для многореференсной награды ту же роль, что VeriFree играет для награды с одним точным ответом. Однако структура градиента меняется: прямое слагаемое теперь распределяется по нескольким эталонным ответам.

Действительно,
\[
    J_{\mathrm{MRVF}}(\theta\mid x,Y^{*})
    =
    \sum_z
    \pi_{\theta}(z\mid x)
    r_{\theta}(x,z,Y^{*}),
\]
поэтому
\begin{align*}
    \nabla_{\theta}J_{\mathrm{MRVF}}
     & =
    \sum_z
    \Big[
        \nabla_{\theta}\pi_{\theta}(z\mid x)
        r_{\theta}(x,z,Y^{*})
        +
        \pi_{\theta}(z\mid x)
        \nabla_{\theta}r_{\theta}(x,z,Y^{*})
        \Big]
    \\
     & =
    \mathbb{E}_{z\sim\pi_{\theta}(z\mid x)}
    \Big[
        r_{\theta}(x,z,Y^{*})
        \nabla_{\theta}\log\pi_{\theta}(z\mid x)
        +
        \nabla_{\theta}r_{\theta}(x,z,Y^{*})
        \Big].
\end{align*}
Теперь раскроем производную по эталонным ответам:
\begin{align*}
    \nabla_{\theta}r_{\theta}(x,z,Y^{*})
     & =
    \nabla_{\theta}
    \sum_{y\in Y^{*}}
    \pi_{\theta}(y\mid x,z)
    \\
     & =
    \sum_{y\in Y^{*}}
    \nabla_{\theta}
    \pi_{\theta}(y\mid x,z)
    \\
     & =
    \sum_{y\in Y^{*}}
    \pi_{\theta}(y\mid x,z)
    \nabla_{\theta}
    \log\pi_{\theta}(y\mid x,z).
\end{align*}
Следовательно,
\[
    \nabla_{\theta}J_{\mathrm{MRVF}}
    =
    \mathbb{E}_{z\sim\pi_{\theta}(z\mid x)}
    \Bigg[
        r_{\theta}(x,z,Y^{*})
        \nabla_{\theta}\log\pi_{\theta}(z\mid x)
        +
        \sum_{y\in Y^{*}}
        \pi_{\theta}(y\mid x,z)
        \nabla_{\theta}\log\pi_{\theta}(y\mid x,z)
        \Bigg].
\]

Если \(r_{\theta}(x,z,Y^{*})>0\), определим веса
\[
    w_{\theta}(y\mid x,z,Y^{*})
    =
    \frac{\pi_{\theta}(y\mid x,z)}
    {r_{\theta}(x,z,Y^{*})}
    \mathbb{I}\{y\in Y^{*}\}.
\]
Тогда \(w_{\theta}\) является распределением по наблюдаемому множеству эталонных ответов, и
\[
    \nabla_{\theta}r_{\theta}(x,z,Y^{*})
    =
    r_{\theta}(x,z,Y^{*})
    \sum_{y\in Y^{*}}
    w_{\theta}(y\mid x,z,Y^{*})
    \nabla_{\theta}\log\pi_{\theta}(y\mid x,z).
\]
Эквивалентно,
\[
    \nabla_{\theta}\log r_{\theta}(x,z,Y^{*})
    =
    \sum_{y\in Y^{*}}
    w_{\theta}(y\mid x,z,Y^{*})
    \nabla_{\theta}\log\pi_{\theta}(y\mid x,z).
\]

Эта форма проясняет отличие от одноэталонного VeriFree. В VeriFree прямое слагаемое --- это градиент \(\log\pi_{\theta}(y^{*}\mid x,z)\) для одного эталонного ответа. В MRVF прямое слагаемое --- это градиент логарифма суммы вероятностей по многим эталонным ответам. Поэтому обновление не отдаёт априорного предпочтения одному эталону; оно увеличивает суммарную вероятностную массу на наблюдаемом множестве, причём каждый эталонный ответ взвешивается текущей вероятностной массой стратегии, назначенной ему.

\subsection{Логарифмическая масса как практическая цель}

Идеальная формулировка MRVF использует вероятности полных последовательностей:
\[
    \pi_{\theta}(y\mid x,z)
    =
    \prod_{t=1}^{L_y}
    \pi_{\theta}(y_t\mid x,z,y_{<t}).
\]
Для коротких проверяемых ответов такая величина может быть удобной. Для шуток она становится проблематичной: полная шутка обычно состоит из десятков токенов, поэтому произведение токеновых вероятностей быстро становится очень малым. Кроме того, ненормированная вероятность последовательности систематически предпочитает более короткие ответы, поскольку каждый дополнительный токен умножает вероятность на число меньше единицы.

Поэтому в практической реализации используется логарифмическая суррогатная цель. Для эталонного ответа \(y_j^{*}\) при запросе \(x\) и траектории рассуждения \(z_i\) сначала вычисляется токеновое логарифмическое правдоподобие:
\[
    \ell_{ij}
    =
    \sum_{t=1}^{L_j}
    \log
    \pi_{\theta}
    \left(
    y_{j,t}^{*}
    \mid
    x,z_i,y_{j,<t}^{*}
    \right).
\]
Чтобы уменьшить смещение в пользу коротких эталонных ответов, используется нормировка по длине:
\[
    s_{ij}
    =
    \frac{1}{L_j^{\gamma}}
    \ell_{ij}.
\]
В основной реализации используется среднее логарифмическое правдоподобие на токен, то есть \(\gamma=1\):
\[
    s_{ij}
    =
    \frac{1}{L_j}
    \sum_{t=1}^{L_j}
    \log
    \pi_{\theta}
    \left(
    y_{j,t}^{*}
    \mid
    x,z_i,y_{j,<t}^{*}
    \right).
\]

Затем для траектории \(z_i\) вычисляется многореференсная логарифмическая масса:
\[
    m_i
    =
    \log
    \sum_{j=1}^{M}
    \alpha_j
    \exp(s_{ij}),
\]
где \(M=|Y^{*}(x)|\), а \(\alpha_j\) --- вес \(j\)-го эталонного ответа. В базовой реализации используются равные веса \(\alpha_j=1/M\). Операция \(\log\sum\exp\) обеспечивает численную устойчивость и позволяет агрегировать несколько эталонных ответов без явного перехода к крайне малым вероятностям полных последовательностей.

Таким образом, идеальная цель MRVF мотивирует суммирование вероятностной массы по множеству эталонных ответов, а реализованная цель использует её устойчивую логарифмическую аппроксимацию. Это изменение не является чисто технической деталью: оно меняет масштаб и геометрию оптимизации. Однако для длинных открытых ответов такая замена необходима, поскольку сырые вероятности последовательностей дают слишком слабый и сильно зависящий от длины сигнал.

\subsection{Неполнота наблюдаемых эталонных ответов}

Наблюдаемое множество \(Y^{*}(x)\) не является полным множеством всех хороших шуток для запроса \(x\). Полезно мысленно различать \(Y^{*}(x)\) и скрытое множество \(Y_{\mathrm{true}}(x)\), содержащее все ответы, которые были бы признаны удачными человеческой аудиторией:
\[
    Y^{*}(x)\subseteq Y_{\mathrm{true}}(x).
\]
Множество \(Y_{\mathrm{true}}(x)\) никогда не наблюдается полностью. Это принципиальное отличие юмора от задач с проверяемым ответом.

Возникают две формы неполноты. Первая --- структурная. Даже если для одного запроса найдено несколько эталонных шуток, существует много других допустимых панчлайнов, отсутствующих в корпусе или не найденных индексом. Вторая --- вычислительная. Даже внутри наблюдаемого множества \(Y^{*}(x)\) на практике может использоваться только ограниченное число эталонных ответов, чтобы уменьшить стоимость принудительного учительского декодирования.

Эта неполнота создаёт риск ложных отрицательных примеров. Оптимизация \(m_i\) или \(r_{\theta}(x,z,Y^{*})\) увеличивает массу на наблюдаемые эталоны. Поскольку распределение модели по всем возможным ответам нормировано,
\[
    \sum_{y\in\mathcal{Y}}
    \pi_{\theta}(y\mid x,z)
    =
    1,
\]
увеличение массы на \(Y^{*}(x)\) может происходить за счёт ответов из \(\mathcal{Y}\setminus Y^{*}(x)\). Но это дополнение содержит не только плохие шутки, но и хорошие, не попавшие в наблюдаемое множество:
\[
    Y_{\mathrm{true}}(x)\setminus Y^{*}(x).
\]
Следовательно, MRVF-обучение по неполному множеству эталонов может непреднамеренно подавлять допустимые, но ненаблюдаемые юмористические продолжения.

Этот риск не является ошибкой вывода градиента. Он является структурным следствием обучения по неполной разметке в области с большим числом допустимых решений. Поэтому результаты MRVF следует интерпретировать не как прямую оптимизацию истинной человеческой смешности, а как оптимизацию наблюдаемой многореференсной суррогатной цели. В экспериментальной части эта проблема учитывается через сравнение с базовыми моделями и автоматическую внешнюю оценку сгенерированных ответов.

\subsection{Использование подмножества эталонных ответов}

Если наблюдаемое множество \(Y^{*}(x)\) велико, вычисление правдоподобия всех эталонных ответов для каждой траектории \(z_i\) может быть дорогим. Поэтому на практике можно использовать подмножество \(Y\subseteq Y^{*}(x)\). Для идеальной вероятностной массы
\[
    r_{\theta}(x,z,Y^{*})
    =
    \sum_{y\in Y^{*}}
    \pi_{\theta}(y\mid x,z)
\]
случайное равномерное подмножество даёт простую связь с полной величиной. Пусть каждый эталонный ответ включается в \(Y\) с вероятностью \(\rho\). Тогда
\[
    r_{\theta}(x,z,Y)
    =
    \sum_{y\in Y^{*}}
    \mathbb{I}\{y\in Y\}
    \pi_{\theta}(y\mid x,z),
\]
и
\[
    \mathbb{E}_{Y}
    \big[
        r_{\theta}(x,z,Y)
        \big]
    =
    \rho
    r_{\theta}(x,z,Y^{*}).
\]
Следовательно, для сырой вероятностной массы случайное подмножество даёт оценку, пропорциональную полной массе в математическом ожидании.

Для логарифмической массы ситуация сложнее. Величина
\[
    m(x,z,Y)
    =
    \log
    \sum_{y\in Y}
    \alpha_y
    \exp(s_{\theta}(y\mid x,z))
\]
не является несмещённой оценкой \(m(x,z,Y^{*})\), поскольку логарифм стоит после суммирования. Поэтому использование подмножества в реализованной цели следует понимать как стохастическое приближение для снижения вычислительной стоимости, а не как строго несмещённый оцениватель полной логарифмической массы. Это ещё одна причина рассматривать MRVF как суррогатную обучающую цель, а не как точную оптимизацию истинного множества всех удачных шуток.

\subsection{Цикл обучения}

Для каждого запроса \(x\) модель выбирает группу из \(G\) траекторий рассуждения:
\[
    z_1,\ldots,z_G
    \sim
    \pi_{\theta}(z\mid x).
\]
Для каждой траектории \(z_i\) вычисляются нормированные логарифмические правдоподобия эталонных ответов \(s_{ij}\), после чего агрегируется многореференсная логарифмическая масса
\[
    m_i
    =
    \log
    \sum_{j=1}^{M}
    \alpha_j
    \exp(s_{ij}).
\]

Слагаемое, отвечающее за выбор траектории рассуждения, имеет вид оценивателя функции оценки. Его задача --- повысить вероятность тех траекторий \(z_i\), для которых наблюдаемое множество эталонных ответов получает большую логарифмическую массу. Для уменьшения зависимости от масштаба наград внутри одного запроса используется групповая нормировка, аналогичная GRPO:
\[
    A_i
    =
    \frac{m_i-\mu_x}{\sigma_x+\varepsilon},
    \qquad
    \mu_x=
    \frac{1}{G}
    \sum_{k=1}^{G}
    m_k,
\]
\[
    \sigma_x
    =
    \sqrt{
        \frac{1}{G}
        \sum_{k=1}^{G}
        (m_k-\mu_x)^2
    }.
\]
Тогда слагаемое для обновления распределения траекторий записывается как
\[
    \mathcal{L}_{z}
    =
    -
    \frac{1}{G}
    \sum_{i=1}^{G}
    \operatorname{sg}(A_i)
    \log
    \pi_{\theta}(z_i\mid x),
\]
где \(\operatorname{sg}(\cdot)\) обозначает остановку градиента.

Остановка градиента здесь принципиальна. Преимущество \(A_i\) используется как коэффициент оценивателя функции оценки, а не как дифференцируемая часть вычислительного графа. Если пропустить градиент через \(A_i\), автоматическое дифференцирование добавит дополнительные слагаемые, не соответствующие исходному выводу. Прямой градиент по эталонным ответам учитывается отдельно, через слагаемое логарифмической массы.

Прямое эталонное слагаемое задаётся как отрицательная логарифмическая масса:
\[
    \mathcal{L}_{\mathrm{ref}}
    =
    -
    \frac{1}{G}
    \sum_{i=1}^{G}
    m_i.
\]
Оно отвечает за увеличение условного правдоподобия эталонных ответов при уже выбранных траекториях рассуждения. В раскрытой форме это слагаемое соответствует градиенту
\[
    \nabla_{\theta}m_i
    =
    \sum_{j=1}^{M}
    \rho_{ij}
    \nabla_{\theta}s_{ij},
\]
где
\[
    \rho_{ij}
    =
    \frac{
        \alpha_j\exp(s_{ij})
    }{
        \sum_{k=1}^{M}
        \alpha_k\exp(s_{ik})
    }.
\]
Таким образом, эталонные ответы внутри \(Y^{*}(x)\) получают веса, пропорциональные их текущему нормированному правдоподобию при данной траектории \(z_i\). Если несколько эталонов имеют сопоставимую массу, градиент распределяется между ними; если один эталон доминирует, обновление концентрируется на нём.

Итоговая функция потерь для одного запроса имеет вид
\[
    \mathcal{L}_{\mathrm{MRVF}}
    =
    \lambda_z
    \mathcal{L}_{z}
    +
    \lambda_{\mathrm{ref}}
    \mathcal{L}_{\mathrm{ref}},
\]
где \(\lambda_z\) и \(\lambda_{\mathrm{ref}}\) управляют относительным вкладом обновления траекторий рассуждения и прямого эталонного слагаемого. В мини-пакете эта величина усредняется по запросам.

Эта формулировка сохраняет основную идею VeriFree: награда получается не от внешней модели-верификатора, а из вероятностей самой обучаемой модели. Отличие состоит в двух изменениях, необходимых для юмора. Во-первых, используется множество эталонных ответов вместо одного. Во-вторых, сырая вероятность полной последовательности заменяется нормированной логарифмической массой, более устойчивой для длинных и вариативных ответов.

\paragraph{Связь с RLOO.}
Альтернативным способом снижения дисперсии является базовый уровень с исключением текущего элемента:
\[
    b_i
    =
    \frac{1}{G-1}
    \sum_{k\ne i}
    m_k,
    \qquad
    A_i^{\mathrm{RLOO}}
    =
    m_i-b_i.
\]
Такой оцениватель ближе к классическому выводу REINFORCE, поскольку вычитает базовый уровень, не зависящий от текущей траектории \(z_i\). Групповая стандартизация, используемая в GRPO,
\[
    A_i
    =
    \frac{m_i-\mu_x}{\sigma_x+\varepsilon},
\]
является более сильной инженерной нормировкой: она не только центрирует, но и масштабирует значения внутри группы. Это помогает стабилизировать обучение при разных запросах, однако меняет точный масштаб оценивателя. В данной работе эта нормировка рассматривается как практическая суррогатная процедура, выбранная для устойчивости обучения.

\subsection{Замечание о базовых уровнях и остановке градиента}

Механизмы снижения дисперсии применимы прежде всего к слагаемому функции оценки по траектории \(z\). Если \(b(x)\) не зависит от выбранной траектории \(z_i\), то
\begin{align*}
    \mathbb{E}_{z_i\sim\pi_{\theta}(\cdot\mid x)}
    \big[
        b(x)
        \nabla_{\theta}
        \log\pi_{\theta}(z_i\mid x)
        \big]
     & =
    b(x)
    \sum_{z_i}
    \pi_{\theta}(z_i\mid x)
    \nabla_{\theta}
    \log\pi_{\theta}(z_i\mid x)
    \\
     & =
    b(x)
    \sum_{z_i}
    \nabla_{\theta}
    \pi_{\theta}(z_i\mid x)
    \\
     & =
    b(x)
    \nabla_{\theta}1
    =
    0.
\end{align*}
Поэтому вычитание базового уровня из награды не меняет математическое ожидание слагаемого функции оценки, но может уменьшить его дисперсию.

Прямое эталонное слагаемое устроено иначе. Оно не является оценивателем по выбранному ответу \(y\sim\pi_{\theta}(\cdot\mid x,z)\), а вычисляется через принудительное учительское декодирование наблюдаемых эталонных ответов. Поэтому наивное вычитание базового уровня из \(\mathcal{L}_{\mathrm{ref}}\) в общем случае вносит смещение.

Например, в случае одного эталонного ответа выражение
\[
    \mathbb{E}_{z}
    \Big[
        (r_{\theta}(x,z,y^{*})-b(x))
        \nabla_{\theta}
        \log\pi_{\theta}(y^{*}\mid x,z)
        \Big]
\]
раскладывается как
\[
    \mathbb{E}_{z}
    \Big[
        r_{\theta}(x,z,y^{*})
        \nabla_{\theta}
        \log\pi_{\theta}(y^{*}\mid x,z)
        \Big]
    -
    b(x)
    \mathbb{E}_{z}
    \Big[
        \nabla_{\theta}
        \log\pi_{\theta}(y^{*}\mid x,z)
        \Big].
\]
Второй член в общем случае не равен нулю. Та же проблема сохраняется в многореференсном случае с градиентом логарифмической массы. Поэтому групповая нормировка используется только в слагаемом \(\mathcal{L}_z\), а прямое эталонное слагаемое вычисляется непосредственно из логарифмической массы.

Та же логика объясняет использование оператора \(\operatorname{sg}\) в \(\mathcal{L}_z\). Значение преимущества используется как числовой коэффициент при \(\log\pi_{\theta}(z_i\mid x)\). Дифференцировать этот коэффициент внутри слагаемого функции оценки не нужно: соответствующий градиент по эталонным ответам уже представлен отдельным слагаемым \(\mathcal{L}_{\mathrm{ref}}\).

В совокупности MRVF наследует ключевое вычислительное преимущество обучения без внешнего верификатора: он избегает внешней модели награды и явной выборки итоговых ответов для проверки. Однако для юмора метод должен быть адаптирован к многореференсному и неполному характеру задачи. Поэтому центральными элементами предлагаемого подхода становятся построение наблюдаемого множества \(Y^{*}(x)\), нормированная логарифмическая масса эталонных ответов и групповая нормировка траекторий рассуждения.